\section{A fortyfold improvement for Suzuki's unsmoothed prime-power race}
\label{sec:suzukimoments}
Write $\Lambda$ for the von Mangoldt function and set
\[
 E(x)=\sum_{n\le x}'\frac{\Lambda(n)}{\sqrt n}-2\sqrt x .
\]
The prime gives half weight to a final prime power. Suzuki asks how often
$E(x)\le0$ in logarithmic measure (\href{https://arxiv.org/html/2411.07436v3}{\S1.1, after (6)}).
The incoming report \texttt{grh\_cm\_four\_terms\_and\_prime\_bias.zip}
proves an upper frequency of $0.00323$ for the opposite sign and a
strictly positive lower frequency. We give an explicit lower bound and reduce the upper bound by using
finite information about additive relations between zero ordinates.

\medskip\noindent\textbf{Theorem.}
Let $p(a,L)=L^{-1}|\{t\in[a,a+L]:E(e^t)>0\}|$ for $a\ge1$. Then
\[
 10^{-937}<\liminf_{L\to\infty}\ \inf_{a\ge1}p(a,L)
 \le\limsup_{L\to\infty}\ \sup_{a\ge1}p(a,L)<0.00007943 .
\]
In particular, the lower logarithmic density of the set $E(x)<0$ exceeds
$99.992057\%$. Its upper logarithmic density is at most $1-10^{-937}$. The existence and exact value of the logarithmic density
at zero remain open in this analysis. The only global zero-location input
is the zeta instance of the Lean theorem. All independence information
used below is finite and verified.

\paragraph{The mean and its Fourier series.}
Write $\gamma$ for a distinct positive ordinate and $m_\gamma$ for its
multiplicity. Suzuki's explicit formula (36) and central value (47) give
\[
 E(e^t)=-c+F(t)+o(1),\qquad
 c=\tfrac12\bigl(\pi/2+\log(8\pi)+\gamma_{\mathrm E}\bigr)>2.686,
 \qquad F(t)=-2\sum_{\gamma>0}\frac{m_\gamma}{\gamma}\sin(\gamma t).
\]
Here $\gamma_{\mathrm E}$ is Euler's constant. This uses the same passage
from zero multiplicities to prime-sum residues that occurs in
\texttt{New/Reduction/Zeros/PrimeError.lean}, with the weight from Suzuki's
question. For clarity, the frequency-series convergence needed below can
be checked directly. For a finite sum $P=\sum A_\lambda e^{i\lambda t}$
put $B_j=\sum_{j\le\lambda<j+1}|A_\lambda|$. Integrating $|P|^2$ against
translates of $(\sin(t/2)/(t/2))^2$, whose Fourier transform is supported
in $[-1,1]$, gives
\[
 \sup_u\int_u^{u+1}|P(t)|^2dt\ll\sum_j B_j^2 .
\]
The kernel is bounded below on $[-1/2,1/2]$ and only equal or neighboring
frequency blocks contribute. The zero count in a unit block is
$O(\log(2+|j|))$, with multiplicities. Hence the symmetric partial sums
converge in this uniform local mean-square norm, with their limit agreeing
almost everywhere with the pointwise explicit formula. Finite-polynomial
means on long intervals converge uniformly in the starting point, by
integrating each exponential. Equivalently, these means are Haar
integrals on the closure of the phase flow
$t\mapsto(e^{i\gamma t})_{\gamma>0}$ in the product of circles. Uniform
mean-square approximation extends this interpretation to the positive
part of any sine-series tail, since $x\mapsto x_+$ is $1$-Lipschitz.
Write $\|\cdot\|_p$ for the Haar norm. Inversion in the phase group
sends every sine sum $Y$ to $-Y$, giving
$\|(Y)_+\|_p=2^{-1/p}\|Y\|_p$ whenever the moment exists.

\paragraph{Moment estimates from finite certificates.}
Number the zeros by increasing positive ordinate. They are simple over
the entire finite range used below. Decompose
$F=P+M_1+M_2+M_3+R$. The indices for $P$ are 1 through 500. Each $M_j$
uses the range in the table, and $R$ is the remaining tail. For a
finite block $I$ put $v_I=2\sum_{i\in I}\gamma_i^{-2}$.

\begin{center}\small
\begin{tabular}{r r r r}
\hline
Indices & $v_I$ (rounded) & Integers sorted & Minimum gap${}^{*}$\\ \hline
1--500 & 0.04391240805 & \multicolumn{2}{l}{Best--Trudgian, Table 2}\\
501--6,000 & 0.00190153507 & 30,250,000 & 2,892\\
6,001--12,000 & 0.00016163380 & 17,997,000 & 236\\
12,001--18,000 & 0.00006286009 & 17,997,000 & 2,478\\ \hline
\end{tabular}

${}^{*}$In units of $10^{-14}$. The arithmetic uses unrounded enclosures.
\end{center}

For the first 500 ordinates, Best and Trudgian's Table 2 excludes every
nontrivial integer relation with coefficients of absolute value at most
$554,000$ (\href{https://arxiv.org/abs/1209.3843}{Math. Comp. 84 (2015), 2433--2446}).
In particular the 18th moment of $P$ equals the moment for independent
uniform phases. Indeed, a monomial in the expanded sine polynomial
has coefficient vector with $\ell^1$ norm at most 18. Only the constant
character integrates to one. For a single independent phase with
amplitude $2/\gamma$, the exponential moment series has coefficients
$1/(\gamma^{2k}(k!)^2)$ in degree $2k$ and zero in odd degree. It is
coefficientwise bounded by $\exp(z^2/\gamma^2)$ as a series in $z$.
Multiplication and degree extraction therefore give
\[
 \mathbb E[P^{2r}]\le (2r-1)!!\,v_P^r,\qquad r=9.
\]

For each middle block we prove that only pair cancellation contributes
to the fourth moment. We computed all 18,000 ordinates as certified balls
with FLINT, retaining the enclosing decimal cell of width $10^{-22}$.
Put $u_i=\lfloor10^{14}\gamma_i\rfloor$, computed from that cell.
The table records the minimum separation among positive differences
$u_j-u_i$ with $j>i$. In the first middle block we include in the same
sort all sums $u_j+u_i$ with $j\ge i$. Different positive differences
separated by more than three units exclude a nontrivial two-against-two
equality between ordinates, since such an equality would give two equal
positive differences after canceling common terms. A three-against-one
equality would give a sum equal to a positive difference. Their scaled
rounding errors are respectively in $[0,2]$ and $[-1,1]$, so the same
gap threshold suffices. In each of the last two blocks $3\min\gamma_i>
\max\gamma_i$ already excludes three-against-one. Four terms with a
common sign have positive absolute sum. Thus the fourth moment of
each $M_j$ equals its independent-phase value and is at most $3v_{M_j}^2$.
This certificate imposes no condition on relations between the blocks.

\paragraph{The part beyond the certificate.}
Brent, Platt and Trudgian give $\sum_{\gamma>0}m_\gamma/\gamma^2\le0.023105$
in equation (20) of \href{https://maths-people.anu.edu.au/brent/pd/rpb274-final.pdf}{J. Number Theory 238 (2022), 740--762}.
Platt and Trudgian verified simplicity to $H=3\cdot10^{12}$
(\href{https://doi.org/10.1112/blms.12460}{Bull. Lond. Math. Soc. 53 (2021)}).
Above $H$ a jump of the counting function is bounded by
$m_\gamma\le0.56\log\gamma$, by Corollary 1 of the 2022 paper.
Its Lemmas 5 and 6 bound the sums of $1/\gamma^2$ and
$\log(\gamma/2\pi)/\gamma^2$ beyond $H$, counting multiplicity.
With $D=\log H$ they give
\[
 \sum_{\gamma>H}\frac{m_\gamma(m_\gamma-1)}{\gamma^2}
 \le\frac{0.56}{2\pi H}\{D^2+(\log(2\pi)-1)D\}<2.524\cdot10^{-11}.
\]
The full variance is $2\sum m_\gamma^2/\gamma^2$ by orthogonality.
Subtracting the four finite block variances, with their certified
endpoints, yields $v_R=\|R\|_2^2<0.000171563042$.

\paragraph{A one-sided estimate without inter-block independence.}
Put $a_0=0.79$, $b_0=0.37$ and $M=\sum_j M_j$. Pointwise,
\[
 (F-a_0-b_0)_+\le(P-a_0)_+ +(M-b_0)_+ + R_+ .
\]
For $h>0$ and $r\ge2$, maximizing $(x-h)/x^r$ for $x\ge h$ gives
$(x-h)_+\le (r-1)^{r-1}(x_+)^r/(r^r h^{r-1})$.
With $r=9$ for $P$ and $r=2$ for $M$, symmetry and Minkowski give
\begin{align*}
 \|(P-a_0)_+\|_2 &\le
 \frac{(r-1)^{r-1}}{r^ra_0^{r-1}}\sqrt{(2r-1)!!\,v_P^r/2}<0.000923214,\\
 \|(M-b_0)_+\|_2 &\le
 \frac{\sqrt{3/2}}{4b_0}\left(\sum_{j=1}^3\sqrt{v_{M_j}}\right)^2<0.003415944,\\
 \|R_+\|_2&=\sqrt{v_R/2}<0.009261832.
\end{align*}
Chebyshev's inequality bounds the requested upper frequency by the
square of the sum of these norms divided by $(c-a_0-b_0)^2$.
To absorb the $o(1)$ in the explicit formula use threshold $c-\eta$,
discard a fixed initial interval at cost $O_\eta(1/L)$, then let
$\eta\downarrow0$. The unrounded computation gives an upper bound
below $0.00007942911$. Equality $E(x)=0$ occurs on a countable set,
since $E$ decreases strictly between consecutive prime powers.

\paragraph{An explicit lower bound for the rare sign.}
Let $G$ be the compact phase group above, with probability Haar measure,
and let $K_T$ be convolution along its one-parameter flow with the
nonnegative real-line Fej\'er kernel of integral one and Fourier multiplier
$(1-|\lambda|/T)_+$. Then $K_T$ preserves nonnegativity and constants,
is self-adjoint, and contracts the supremum norm. For $T=1420$ the first
$1{,}000$ ordinates lie below the cutoff, with the next above. Put
$d_i=2(1/\gamma_i-1/T)$ for $1\le i\le 1{,}000$. The smoothed sum
$K_TF$ is the finite polynomial $-\sum d_i\sin\theta_i$, where
$e^{i\theta_i}$ is coordinate $i$ on $G$.

Set $k=23$, and normalize $w(\theta)=(1-\sin\theta)^k$ to have
mean one on the circle. Direct integration gives
\[
  \max w=4^k/\binom{2k}{k}=:B,\qquad
  \int w(\theta)(-\sin\theta)\,\frac{d\theta}{2\pi}=k/(k+1).
\]
Indeed the circle mean before normalization is
$2^{-k}\binom{2k}{k}$, and the second identity follows by taking the
ratio of the successive means and subtracting one. Set
$W=\prod_{i=1}^{500}w(\theta_i)$. The finite Best--Trudgian certificate
applies coordinate by coordinate to every monomial of $W$ and of
$W\sin\theta_i$ for $i\le500$, since no exponent exceeds 24. Thus
$\int_G W=1$, while $0\le W\le B^{500}$. For the remaining 500 modes we
use only $|\sin\theta_i|\le1$. Therefore
\[
 \int_G W(K_TF-c)\ge
 \Delta :=\frac{k}{k+1}\sum_{i=1}^{500}d_i
     -\sum_{i=501}^{1000}d_i-c>0.00091075.
\]
Moving the convolution onto $W$ and using $0\le K_TW\le B^{500}$ gives
\[
 \Delta\le B^{500}\int_{F>c}(F-c)\le
 B^{500}\|F\|_2\,\mathbb P(F>c)^{1/2}.
\]
The arithmetic verifies $\Delta^2/(B^{1000}\|F\|_2^2)>10^{-937}$ using
the preceding upper bound for $\|F\|_2^2$.
To pass to uniform lower frequency, bounded Lipschitz test functions of
$F$ have Haar means uniformly on long intervals by the local mean-square
approximation above. Approximate the indicator of the open half-line
$(c,\infty)$ from below by such functions. The decaying explicit-formula
remainder also tends to zero in mean, uniformly on long intervals
starting at $a\ge1$. This proves the lower half of the theorem.
